\documentclass[12pt, a4paper]{article}
\usepackage[utf8]{inputenc}
\usepackage{amsmath}
\usepackage{graphicx}
\usepackage{hyperref}
\usepackage{geometry}
\geometry{a4paper, margin=1in}
\usepackage{titlesec}
\usepackage{xcolor}
\usepackage{listings}

\title{\textbf{Packet Filtering Firewall: Phase 1 Architectural and Technical Analysis}}
\author{Bikesh}
\date{\today}

\begin{document}

\maketitle

\begin{abstract}
This report provides a technical analysis of the architecture, design, and internal implementation of a custom packet filtering firewall (Phase 1). The system operates on a decoupled architecture, utilizing a Python-based \texttt{FastAPI} and \texttt{Scapy} backend for network-layer packet inspection, alongside a React 18 frontend built upon IBM's Carbon Design System. The application integrates stateless policy enforcement, threshold-based rate limiting, data persistence, modular state management, and an administrative visualization interface.
\end{abstract}

\tableofcontents
\newpage

\section{Introduction}
This project implements a real-time packet filtering firewall system. By decoupling the network-layer packet inspection engine from the web-based administrative console, the firewall allows network administrators to manage, test, and visualize network traffic via a user interface without blocking the performance of the core packet-processing daemon. This document outlines the Phase 1 implementation.

\section{Backend Inspection Engine (Python/FastAPI)}
The backend operates at the network layer and is decoupled into discrete modules handling packet processing, database transactions, and API endpoints.

\subsection{Packet Sniffing and Processing (\texttt{engine.py})}
Built upon the Python \texttt{Scapy} library, the packet processing engine captures inbound and outbound network traffic inside a background thread.
\begin{itemize}
    \item \textbf{Packet Extraction:} Extracts Layer 3 and Layer 4 attributes, standardizing IPs, protocols (TCP, UDP, ICMP), and ports into a dictionary structure.
    \item \textbf{Rate Limiting Mechanism:} An in-memory sliding window algorithm tracks packet counts per IP address. The engine differentiates between two thresholds: \texttt{flood\_threshold} (per second limits) which triggers an automatic IP addition to the Blocklist, and \texttt{rate\_limit} (per minute limits) which drops subsequent traffic.
    \item \textbf{Rule Evaluation:} Packets are processed sequentially against a cache of active firewall rules. The engine executes the first matching disposition (\texttt{ALLOW}, \texttt{DROP}, \texttt{BLOCK}) and logs the outcome.
\end{itemize}

\subsection{Database Management (\texttt{database.py})}
Data persistence is handled by \texttt{aiosqlite}, ensuring non-blocking I/O database transactions.
\begin{itemize}
    \item \textbf{Rules Table:} Stores active firewall policies enforcing protocol, source/destination IP, and port constraints.
    \item \textbf{Blocklist Table:} Maintains a ledger of blocked IPs, utilizing UUIDs for primary keys.
    \item \textbf{Logs Table:} A ledger of processed packets including their evaluation outcome.
    \item \textbf{Settings Table:} A Key-Value store managing engine-level configurations.
\end{itemize}

\subsection{Data Validation (\texttt{models.py})}
The backend employs Pydantic schemas to validate JSON payloads. It enforces camelCase naming conventions across API boundaries (e.g., \texttt{srcIp}, \texttt{dstPort}) to maintain consistency between the frontend and backend.

\subsection{REST API Routers (\texttt{api/})}
The backend exposes API endpoints via FastAPI routers, secured via JSON Web Tokens (JWT).
\begin{itemize}
    \item \textbf{\texttt{/capture}}: Handles engine telemetry (\texttt{/status}, \texttt{/stats}, \texttt{/packets}) and state toggling (\texttt{/start}, \texttt{/stop}, \texttt{/clear}).
    \item \textbf{\texttt{/test}}: Allows the frontend to POST simulated packets to test firewall disposition logic.
    \item \textbf{\texttt{/rules}, \texttt{/blocklist}, \texttt{/logs}, \texttt{/settings}}: Provide CRUD capabilities. Configuration updates require atomic POST operations per setting.
\end{itemize}

\section{Frontend Web Interface (React/Carbon)}
The administrative frontend is a single-page application built on React 18, compiled via Vite, and styled with the IBM Carbon Design System.

\subsection{State Management (\texttt{context/})}
The architecture uses the React Context API to manage global state.
\begin{itemize}
    \item \textbf{\texttt{AuthContext.jsx}:} Governs the JWT \texttt{fw\_token} lifecycle and provides an \texttt{authFetch} interceptor utility that injects \texttt{Bearer} authorization headers into outgoing requests.
    \item \textbf{\texttt{AppContext.jsx}:} Acts as the centralized telemetry hub. It initializes polling intervals (e.g., fetching graphs every 1000ms), maintains the traffic history array, and caches global datasets such as Rules and Blocklists.
\end{itemize}

\subsection{Application Shell \& Layout (\texttt{components/})}
The \texttt{AppShell} manages the responsive layout. It dynamically controls Carbon's \texttt{<SideNav>} utilizing the \texttt{isRail} property, ensuring navigational icons remain visible on the far-left rail when the drawer is collapsed.

\subsection{Feature Modules (\texttt{pages/})}
\begin{itemize}
    \item \textbf{Dashboard:} Subscribes to the \texttt{AppContext} stats object to map traffic throughput into data tables.
    \item \textbf{Live Capture:} Presents a table of packet traffic currently traversing the engine. Implements unified tag coloring: \texttt{ALLOW} (Green), \texttt{BLOCK} (Red), and \texttt{DROP} (Magenta). Invokes the \texttt{/capture/clear} API route to flush the live buffer.
    \item \textbf{Packet Tester:} A simulation tool where administrators construct a packet and query the backend \texttt{/test} endpoint to debug rule logic.
    \item \textbf{Rules \& Blocklists:} Interactive forms to provision policies. The components translate internal state into the camelCase JSON schemas expected by the backend Pydantic models. Deletions map to database UUIDs.
    \item \textbf{Settings:} Orchestrates configuration adjustments by looping through modified input states and dispatching individual HTTP POST events per atomic setting.
\end{itemize}

\section{System Integration Analysis}
A focus of the Phase 1 implementation involved standardizing data transmission between the decoupled subsystems.
\begin{itemize}
    \item \textbf{Payload Synchronization:} React components format outgoing data to match backend expectations, specifically translating snake\_case to camelCase parameters and utilizing UUIDs for targeted deletions.
    \item \textbf{Buffer Polling:} To handle the volume of kernel-level packet sniffing, the backend aggregates captured traffic into discrete chunks before exposing them via the \texttt{/capture/stats} endpoint. The frontend polls this endpoint to calculate network throughput differentials.
    \item \textbf{State Consistency:} Visual statuses across the Dashboard, Live Capture, Logs, and Rules utilize identical mapping functions, ensuring uniform data representation.
\end{itemize}

\section{Conclusion}
The Phase 1 Packet Filtering Firewall project establishes a functional, decoupled administration system. It securely divides the Scapy packet inspection engine from the React-based WebUI, offering a baseline platform for traffic monitoring and rule enforcement. Future iterations (Phase 2 and beyond) will explore advanced features such as Deep Packet Inspection (DPI) for Layer 7 protocol analysis and integrated anomaly detection.

\end{document}
